% Source: https://artofproblemsolving.com/community/c3h3248919_2024_emcc_individual_speed_test__exeter_math_club_competition
% Stable identity: aops-topic-3248919-q20
Because $B,M,C$ are collinear, $\angle AMD=180^\circ-24^\circ-66^\circ=90^\circ$. Let $N$ be the midpoint of $AD$. In right triangle $AMD$, the midpoint of its hypotenuse is equidistant from all three vertices, so $NA=NM$. Also $NM\parallel AB$, by the trapezoid midsegment theorem. Hence $\angle NAM=\angle AMN=\angle BAM$; write this common angle as $\alpha$.
The angle at $A$ of the trapezoid is $2\alpha$, and by its isosceles symmetry $\angle ABC=2\alpha$. In triangle $ABM$, therefore, $\alpha+2\alpha+24^\circ=180^\circ$, giving $\alpha=52^\circ$. Thus $\boxed{\angle ABC=104^\circ}$.
\begin{asy}
usepackage("amsmath");
size(7cm);
real t=tan(76*pi/180),u=(2*t*t+10)/(t*t-3),d=(u+sqrt(u*u-4))/2,h=(d-1)*t;pair A=(-1,h),B=(1,h),C=(d,0),D=(-d,0),M=(B+C)/2,Np=(A+D)/2;draw(A--B--C--D--cycle);draw(A--M--D,blue);draw(Np--M,dashed);dot(M);dot(Np);label("$A$",A,NW);label("$B$",B,NE);label("$C$",C,SE);label("$D$",D,SW);label("$M$",M,E);label("$N$",Np,W);label("$24^\circ$",M+(-0.35,0.7));label("$66^\circ$",M+(-0.45,-0.65));
\end{asy}
